% Part: computability
% Chapter: computability-theory
% Section: complete-ce-sets

\documentclass[../../../include/open-logic-section]{subfiles}

\begin{document}

\olfileid[fr]{cmp}{thy}{cce}
\olsection{Ensembles effectivement énumérables complets}

\begin{defn}
Un ensemble~$A$ est dit \emph{effectivement
énumérable complet} pour la réductibilité
many-one si
\begin{enumerate}
\item $A$ est effectivement énumérable;
\item pour tout autre ensemble effectivement
  énumérable~$B$, on a $B\leq_m A$.
\end{enumerate}
\end{defn}

Autrement dit, les ensembles effectivement
énumérables complets sont les plus « difficiles »
de cette classe. Ils permettent de répondre aux
questions concernant \emph{n'importe quel}
ensemble effectivement énumérable.

\begin{thm}
Les ensembles $K$, $K_0$ et~$K_1$ sont tous
effectivement énumérables complets.
\end{thm}

\begin{proof}
Pour montrer que $K_0$ est complet, soit $B$
un ensemble effectivement énumérable quelconque.
Il existe alors un indice~$e$ tel que
\[
B = W_e = \Setabs{x}{\cfind{e}(x) \fdefined}.
\]
Posons $f(x)=\tuple{e,x}$. Pour tout~$x$,
on a $x\in B$ si et seulement si $f(x)\in K_0$.
Ainsi $f$ réduit $B$ à~$K_0$.

Dans la preuve de la \olref[k1]{prop:k1}, on réduit
$K_0$ à~$K_1$. La transitivité
(\olref[ppr]{prop:trans-red}) montre donc que
tout ensemble effectivement énumérable se
réduit à~$K_1$, qui est lui-même effectivement
énumérable par cette proposition.

Pour $K$, il faut réduire $K_0$ à~$K$ dans ce
sens. Considérons la fonction calculable partielle
ternaire $q(e,x,z)\simeq\cfind{e}(x)$, qui ignore
$z$, et choisissons un indice~$u$ de~$q$.
Le théorème $s$-$m$-$n$ donne une fonction
totale calculable $r(e,x)=s^2_1(u,e,x)$ telle que,
pour tout~$z$,
$\cfind{r(e,x)}(z)\simeq\cfind{e}(x)$.
En prenant $z=r(e,x)$, on obtient
$\tuple{e,x}\in K_0$ si et seulement si
$r(e,x)\in K$. Fixons aussi un indice de la
fonction nulle part définie. La fonction qui
associe $r((w)_0,(w)_1)$ à chaque code de
paire~$w$ et cet indice à tout autre entier
réduit donc $K_0$ à~$K$. Comme $K$ est
effectivement énumérable, cela prouve sa
complétude.\footnote{Correction éditoriale : la source invoque ici la réduction de $K$ à $K_0$, qui ne prouve pas la complétude de $K$. Nous construisons la réduction nécessaire de $K_0$ à $K$, en traitant aussi les entiers qui ne codent pas une paire.}
\end{proof}

\begin{prob}
Donner une réduction de $K$ à~$K_0$.
\end{prob}

\begin{digress}
Tous les exemples d'ensembles effectivement
énumérables rencontrés jusqu'ici sont donc soit
calculables, soit complets. C'est surprenant :
existe-t-il des ensembles effectivement
énumérables qui ne sont ni calculables ni
complets pour la réduction many-one ? Oui : les ensembles simples de Post en fournissent des exemples.
Le résultat obtenu au milieu des années 1950, indépendamment
par Friedberg et Muchnik, concerne la non-complétude pour la réduction de Turing.\footnote{Correction éditoriale : la source confond les deux notions de complétude. Sur la non-complétude des ensembles simples pour des réductions même plus générales que many-one, voir An.~A.~Muchnik et S.~E.~Positsel'skii, « A class of enumerable sets », \emph{Russian Mathematical Surveys} 54(3), 1999, p.~640, théorème~1, \url{https://doi.org/10.1070/rm1999v054n03ABEH000163}.}
\end{digress}

\end{document}
